\documentclass[12pt]{article}
\usepackage[T1]{fontenc}
\usepackage[utf8]{inputenc}
\usepackage{lmodern}
\usepackage{textcomp}
\usepackage{enumitem}
\usepackage{geometry}
\usepackage{graphicx}
\usepackage{amsmath, amssymb}
\geometry{margin=1in}
\begin{document}
\begin{center}
Political Science - Graduate Level Assessment 24 \
Total Marks: 40 \
Answer all questions.
\end{center}

\section*{Multiple Choice (10 marks)}
\begin{enumerate}[label=\arabic*.]
\item What fact impinged upon America's 'neutral' stance towards the belligerent in World War I between 1915 and 1917?
\begin{enumerate}[label=(\alph*)]
    \item US banks funded the allies much more than Germany and Austria-Hungary
    \item US banks funded Germany and Austria-Hungary much more than the allies
    \item Wilson secretly favoured the German position
    \item None of the above
\end{enumerate}
\item What is meant by the phrase 'empire by invitation'?
\begin{enumerate}[label=(\alph*)]
    \item Voluntary reliance on an external power for security
    \item Willful openness to colonization
    \item Cultural imperialism
    \item Open advocacy of imperialism for economic gain
\end{enumerate}
\item According to realists, what is the fundamental difference between the international system and the domestic system?
\begin{enumerate}[label=(\alph*)]
    \item Armed conflict
    \item Anarchy
    \item Institutions
    \item No common language
\end{enumerate}
\item Peace, commerce, and honest friendship with all nations, entangling alliances with none'. Identify the speaker.
\begin{enumerate}[label=(\alph*)]
    \item James Madison
    \item Abraham Lincoln
    \item Woodrow Wilson
    \item Thomas Jefferson
\end{enumerate}
\item Within American politics, the power to accord official recognition to other countries belongs to
\begin{enumerate}[label=(\alph*)]
    \item the Senate.
    \item the president.
    \item the Secretary of State.
    \item the chairman of the Joint Chiefs.
\end{enumerate}
\end{enumerate}

\section*{True / False (10 marks)}
\textit{Answer True or False}
\begin{enumerate}[label=\arabic*.]
\setcounter{enumi}{5}
\item True or False: In the Vietnam War, technological superiority was the primary factor in the rapid defeat of North Vietnamese forces.
\item True or False: The term 'indiscriminate attack' is generally viewed as undermining the legitimacy of military actions in the context of RMA.
\item True or False: The US consistently ignored the political dynamics and needs of Third World nations during the Cold War era.
\item True or False: Systemic theories assert that a state's foreign policy is primarily shaped by the character of its leader.
\item True or False: The cyber-security discourse emerged in the 1970 s as politicians and academics began to address its significance.
\end{enumerate}

\section*{Long Form Response (20 marks)}
\textit{Answer each question with a clear and complete explanation.}
\begin{enumerate}[label=\arabic*.]
\setcounter{enumi}{10}
\item Discuss the structure of the United Nations Security Council, focusing on the roles and powers of its permanent and rotating members, and analyze the implications of this structure on global governance.
\item Discuss the significance of new critical schools of security studies in shaping the field. Have they integrated into International Relations, or do they retain a distinct identity within security studies?
\end{enumerate}

\vfill
\noindent\textit{End of Paper}
\end{document}